At a tame corank-one wall the cometric on covectors acquires a one-dimensional radical $\mathcal{L}$, the Clifford sheaf continues across the wall, the physical metric has a simple pole along $\mathcal{L}\otimes\mathcal{L}$ with $\mathcal{L}^{\otimes 2} \cong N^*\Sigma$, and the completion adjoins the Witt partner $k^*$ of the radical covector $k$ exactly there, with $kk^* + k^*k = 1$, that is $\ip{k}{k^*} = \tfrac12$ (the ambient basis of Proposition \claimref{dark-energy.sheaf-papers-wall-soldering-variables}(iii) rescales $k^*$ to $\ip{k}{k^*} = 1$, which changes nothing); the transverse sphere is carried by the bilinear matrix algebra of the boundary modes through the principal embedding, so that every angular operator of the lift is a boundary-mode bilinear under the adjoint $\mathfrak{su}(2)$; this last clause, the sphere theorem, is derived in the entropy paper under its Indifference principle (\cite[Principle~5.4]{ButlerT8}: a construction may not carry distinctions that nothing in the framework determines), which that paper declares a premise and not a theorem, and it is conditional on that premise wherever it is used here. \emph{Provenance:} \cite[Thm.~3.1, Cor.~4.5, Thm.~6.2]{ButlerT4}; \cite[Prop.~5.2, Principle~5.4, Thm.~5.5, Prop.~5.6]{ButlerT8}. (The degeneration of the soldering form used in Theorem \claimref{dark-energy.walls-free-sector} is the dual type, a metric on tangent vectors with a radical; Proposition \claimref{dark-energy.sheaf-papers-wall-soldering-variables} states the relation.)
